
\section{Appendix}

\subsection{On Admission to the Congregation}

[At the Free Church’s annual meeting in 1898, which among other things discussed a draft for a congregational constitution, there was discussion of what conditions a congregation had the right to require of those who applied for admission to the congregation. Some held that the congregation ought to require a confession, for example: “I believe in Jesus,” “I know that I love Jesus,” or the like. In the course of the conversation Sverdrup expressed himself as follows, according to the report in Foldebladet for June 22, 1898:]

“One may also speak falsely about Jesus. One may adopt certain over-bold ways of speaking about Jesus and mean little by them; so neither is that any sure sign of true life in God.

The difficulty is not admission to the congregation, but the difficulties must be placed in connection with the congregation itself. Who is to admit? What is the congregation like?

When the congregation is worldly and the pastor is worldly, then they naturally admit such as are like themselves. If, on the other hand, one had an assembly of believers and a believing pastor, there would not be such great difficulties with admission to the congregation.

As conditions are in most places, it is simply impossible to get matters into proper order. An abuse is abroad. Many enter the congregation in order to obtain reassurance over against eternity; others want protection against attacks upon sins and vices within the congregation. When, for example, a pastor rises up against acts of moral recklessness, then people are brought into the congregation in order to drive the pastor away.

It is therefore right to be very careful about admission to the congregation; above all, that one admit man by man, person by person. Admission in this connection means nothing other than that one gives those who are admitted occasion to be among those who govern the congregation, not that one gives them an endorsement as Christians or gives them access to hear God’s Word.

The stricter one is, the better. One thereby hinders no one from salvation and blessedness.

One can be saved even if one does not have one’s name in the congregational register.

According to God’s Word we see that those who were baptized and saved were also added to the congregation. If now some congregation wished to draw a boundary and say: only when you have come to conscious faith may you take part in governing the congregation, then it surely has the right to do so.

I should prefer to write the point differently, so that it might come into agreement with VII, 4, where it says that the congregational council shall ‘examine whether those who request admission to the congregation have the requisite conditions for it.’ There ought to be some testimony, and preferably the testimony of those who have tested the person concerned.”

---

[In accordance with what has been communicated above, Sverdrup at the Free Church’s 9th annual meeting, 1905, gave an address on the conditions for the formation of a congregation, in which he says:]

“It lies in the nature of the matter that, as the congregations grew older and increased in number and esteem, and there came to be many different reasons that might move a person to join the congregation, it had to seek to guard itself against the danger that lies in admitting such members as are drawn to the congregation not by God’s Spirit, but by fleshly reasons. Everything that lives and grows has within itself the instinct of self-preservation and seeks to keep away and to remove what is called ‘foreign matter.’ It is therefore also self-evident that the highest and most delicate of all organisms, Christ’s body, or the congregation, was careful both with admission into and exclusion from the congregation. And partly through its office-bearers and partly through the congregation’s meetings, the congregation exercised the oversight which it recognized as its holy duty toward the Lord and His body. But this did not, of course, conflict with the principle of voluntariness; it only imposed a necessary and justified test upon the character of that voluntariness, whether it was the work of the Spirit or of the flesh.”
